%%
%% Chapter 3: Method -- the GT-Conditioned CVAE algorithm
%% Independent file, included by ../main.tex
%%

\section{Method: The GT-Conditioned CVAE}
\label{sec:method}

We first fix notation and review the deterministic splatting backbone
(Section~\ref{sec:method:backbone}), then present the probabilistic
neural-process formulation (Section~\ref{sec:method:prob}), the
GT-conditioned posterior (Section~\ref{sec:method:p5}), the training
objective (Section~\ref{sec:method:loss}), and the deployment-time
estimator (Section~\ref{sec:method:inference}).

\subsection{Problem Setting and Notation}
\label{sec:method:problem}

Let $x\in\mathbb{R}^{D\times H\times W}$ denote the target CT volume and
let $C=\{(\mathcal{I}_i,\theta_i)\}_{i=1}^{m}$ denote the observed
\emph{context}: $m$ projection images $\mathcal{I}_i$ together with
their gantry angles $\theta_i$. In our protocol $m=6$ fixed views with
raw IDs $\{0,8,16,25,33,41\}$. During training an additional
\emph{dense} set $D$ of $48$ views is available, and the ground-truth
volume $Y$ itself may be encoded by the posterior network. We write
$p_{\mathrm{data}}(x\mid C)$ for the target
posterior we wish to approximate.

\subsection{Deterministic Splatting Backbone}
\label{sec:method:backbone}

We adopt the DIF-Gaussian architecture \cite{lin2024difgaussian} as the
reconstruction backbone. A multi-view image encoder $E$ (\emph{UNet} with
output channels 128) extracts per-view convolutional features; a
geometry-aware multiview transformer with absolute-angle positional
embedding and a spatial window of $\pm1$ anchor neighbors fuses the view
features into a set of $K=512$ anchor tokens at volume resolution 12;
after anchor-level fusion the tokens are unprojected into 3D Gaussian
primitives (position, opacity, anisotropic covariance,
spherical-harmonic radiance). A point decoder consumes per-Gaussian
features together with query point coordinates and predicts per-point
density/attenuation, which is normalized by a \emph{volume rendering}
step to produce the final volume. For a fixed sparse view set, the
deterministic model computes

\begin{equation}
    \hat{x} = g\big(f_{\mathrm{GS}}(E(C))\big),
    \label{eq:deterministic}
\end{equation}

where $f_{\mathrm{GS}}$ denotes the Gaussian formation and $g$ the
point-decoder plus rendering. The backbone is trained by
volume-space mean-squared error (MSE) against $x$.

\subsection{Probabilistic Formulation: Latent Neural Process}
\label{sec:method:prob}

To obtain reconstruction uncertainty we introduce a global latent
variable $z\in\mathbb{R}^{d_z}$ ($d_z=64$) and factorize the model as a
latent neural process \cite{garnelo2018neural,garnelo2018latent} over
the view sets:

\begin{align}
    p(x, z \mid C) &= p(z\mid C)\, p(x\mid z, C),
    \label{eq:joint}\\
    p(z\mid C) &= \mathcal{N}\big(z; \mu_\phi(\psi(C)),\,
        \mathrm{diag}\,\sigma_\phi^2(\psi(C))\big),
    \label{eq:prior}
\end{align}

where $\psi(C)=\big[\,\mathrm{mean}_{k}\, a_k(C);\ \mathrm{obs}(C)\,\big]$
concatenates the pooled anchor features with a 7-dimensional
permutation-invariant observation descriptor
$\mathrm{obs}(C)=\big[\log(1{+}|C|),\, \lambda_{\mathrm{cov}},\,
\sin\bar\theta, \,\cos\bar\theta,\, \sin 2\bar\theta,\,
\cos 2\bar\theta,\, r\big]$ --
count, circular coverage, first/second-order angle moments, and mean
resultant length -- and $\mu_\phi,\sigma_\phi$ are the \emph{prior head},
a bounded Gaussian head with per-dimension sigmoid-bounded
$\sigma\in[\sigma_{\min},\sigma_{\max}]=[0.001,3.0]$ and
$\sigma_{\mathrm{init}}=0.5$. Conditioning the decoder on $z$ is done by
feature-wise linear modulation (FiLM): a modulation network with
zero-mean initialization at scale $10^{-3}$ transforms the Gaussian
feature vector into per-channel scale/shift applied to
$f_{\mathrm{GS}}(E(C))$ before volume rendering:
$f_{\mathrm{GS}}^{(z)}(C)=\mathrm{FiLM}(z)\,f_{\mathrm{GS}}(E(C))$. The
rendered predictive distribution is Gaussian with a fixed observation
variance $\sigma_{\mathrm{obs}}^2=(0.1)^2$.

\subsection{GT-Conditioned CVAE Posterior}
\label{sec:method:p5}

Training with the prior alone ($z\sim p(z\mid C)$) learns a stochastic
mapping from sparse views but gives the model no hint about the
\emph{target}. The conditional variational principle fixes this: we
train a separate \emph{posterior head} so that $z$ is drawn from
$q(z\mid C, Y)$ instead, concentrating $z$ on atypical target content
and maximizing the ELBO. The \emph{GT-conditioned} CVAE posterior
is

\begin{equation}
    q(z\mid C, Y) = \mathcal{N}\big(z; \mu_\phi'(\psi(C)\oplus h_Y), \,
        \mathrm{diag}\,\sigma_{\phi'}^2(\psi(C)\oplus h_Y)\big),
    \label{eq:posterior}
\end{equation}

where $h_Y=\mathrm{GTEncoder}(Y)\in\mathbb{R}^{32}$ is a global
descriptor produced by a lightweight 3D CNN ($1\to16\to16\to32$ feature
channels, two stride-2 downsamplings, GroupNorm, SiLU) that mean-pools
its feature volume over the spatial axes. During training the posterior
head receives the concatenation $\psi(C)\oplus h_Y$ and the model is
optimized with latent samples from $q$; the sharing of the sparse
representation $\psi(C)$ between prior and posterior anchors the two
distributions to the same context, which keeps the KL term bounded and
makes $q$ strictly more informative than $p$.

Critically, the GT-conditioned posterior consumes the ground-truth
volume \emph{only inside the posterior head and only during training}.
The decoder never sees $h_Y$; it observes $z$ and, by default, the
features of the context set. At deployment the GT encoder is not
instantiated and $z$ is sampled from $p(z\mid C)$
(Section~\ref{sec:method:inference}), so inference requires exactly the
same inputs as the deterministic backbone: six sparse views and nothing
else.

\subsection{Training Objective}
\label{sec:method:loss}

The full objective combines, per training step,
\begin{itemize}
\item the \emph{posterior reconstruction} term
      $\mathcal{L}_{\mathrm{vol}}^{(q)}=\mathbb{E}_{q(z\mid C,Y)}\,
      \big[\tfrac{1}{2\sigma_{\mathrm{obs}}^2}\,\|g(f_{\mathrm{GS}}^{(z)}(C))-x\|_2^2\big]$,
      evaluated with one reparameterized sample ($S=1$) over $1024$
      query points of a random volume block;
\item a \emph{weighted KL} term with a linear warm-up,
      $\beta_t = \beta_f\cdot\min(1, t/500)$, $\beta_f = 2\times10^{-2}$,
      where the per-dimension KL is floored by a free-bits allowance of
      $0.5$ nats \cite{kingma2016freebits} before being summed over the
      latent dimensions and averaged over the batch;
\item an optional \emph{prior-reconstruction} term that supervises the
      deployment path by rendering \emph{prior} samples
      $z\sim p(z\mid C)$ against the same target, weighted by
      $w_{\mathrm{prior}}$.
\end{itemize}
Formally,

\begin{equation}
    \mathcal{L}
    = \underbrace{\mathcal{L}_{\mathrm{vol}}^{(q)}}_{\text{posterior rec.}}
    + \beta_t\,
      \underbrace{\mathrm{KL}\big[q(z\mid C,Y)\,\|\,p(z\mid C)\big]}_{\text{KL}}
    + w_{\mathrm{prior}}\,
      \underbrace{\mathcal{L}_{\mathrm{vol}}^{(p)}}_{\text{prior rec.}},
    \label{eq:loss}
\end{equation}

with the KL computed in closed form between diagonal Gaussians. In the
configuration studied here $\beta_f=2{\times}10^{-2}$ and
$w_{\mathrm{prior}}=0$ by default. All parameters
are trained jointly with AdamW ($\mathrm{lr}=2{\times}10^{-4}$, weight
decay $10^{-2}$ for the shared backbone, none for fusion modules) and a
per-epoch StepLR schedule that decays the learning rate by
$(10^{-3})^{1/400}$ per epoch.

\subsection{Deployment-Time Estimator}
\label{sec:method:inference}

At inference the posterior head, the GT encoder, and the dense views are
all absent. We draw $N=8$ Monte-Carlo samples $\{z^{(n)}\}_{n=1}^{N}
\overset{\text{i.i.d.}}{\sim} p(z\mid C)$ from the prior head, render
$N$ volumes, and report the mean prediction and the per-point predictive
variance

\begin{equation}
    \hat{x}
    = \tfrac{1}{N}\textstyle\sum_{n} \hat{x}^{(n)},
    \qquad
    \widehat{\mathrm{Var}}(x)
    = \tfrac{1}{N-1}\textstyle\sum_{n}
      \big(\hat{x}^{(n)}-\hat{x}\big)^2 + \sigma_{\mathrm{obs}}^2 ,
    \label{eq:estimator}
\end{equation}

where the second term is the observation noise.

Algorithm~\ref{alg:p5} summarizes the training and inference
procedure.

\begin{algorithm}[t]
\caption{GT-Conditioned CVAE Training and Inference}
\label{alg:p5}
\begin{algorithmic}[1]
\REQUIRE Sparse context $C$ (6 views); training only: dense set $D$ (48
views), GT volume $Y$; backbone $g\circ f_{\mathrm{GS}}$, heads
$\mu_\phi,\sigma_\phi,\mu_{\phi'},\sigma_{\phi'}$, GT encoder $E_Y$.
\STATE {\bf Training step:}
\STATE \quad $\psi(C)=\big[\mathrm{mean}_k a_k(C);\,\mathrm{obs}(C)\big]$;\;
$h_Y = E_Y(Y)$
\STATE \quad draw $z\sim q(z\mid C,Y)$ by reparameterization of
Eq.~\eqref{eq:posterior}
\STATE \quad decode with FiLM-modulated context features:
$\hat{x}^{(q)}=g(f_{\mathrm{GS}}^{(z)}(C))$;
\STATE \quad (optional) draw $z_p\sim p(z\mid C)$ and
$\hat{x}^{(p)}=g(f_{\mathrm{GS}}^{(z_p)}(C))$
\STATE \quad update all parameters with
$\nabla\,\mathcal{L}$~(\ref{eq:loss})
\STATE {\bf Inference:}
\STATE \quad discard $E_Y,q$; draw $z^{(1)}\dots z^{(N)}\sim p(z\mid C)$
\STATE \quad $\hat{x}=\tfrac1N\sum_n g(f_{\mathrm{GS}}^{(z^{(n)})}(C))$;
estimate variance per Eq.~\eqref{eq:estimator}
\RETURN $\hat{x}$ and per-point variance
\end{algorithmic}
\end{algorithm}

Figure~\ref{fig:p5_flow} summarizes the complete information flow of the
method at training and at deployment, highlighting which modules
consume the ground truth and which ones are discarded before inference.

\begin{figure}[H]
\centering
\caption{Information flow of the GT-conditioned CVAE.
\emph{Top (Training):} the GT encoder and
the posterior head are the only modules that consume the ground-truth
volume $Y$; the decoder always renders from the sparse context
features, so $z$ is the \emph{sole} channel of dense/target
information. The loss combines posterior reconstruction with the
annealed, free-bits-floored $\mathrm{KL}[q\|p]$ and an optional
prior-reconstruction term. \emph{Bottom (Deployment):} the GT encoder
and posterior are removed; $N{=}8$ samples $z_n\sim p(z\mid C)$ from
the same sparse-view prior yield a mean volume and per-point variance.}
\label{fig:p5_flow}
\resizebox{\textwidth}{!}{%
\begin{tikzpicture}[
    font=\small,
    >=stealth,
    box/.style={draw, rounded corners=2pt, align=center, inner sep=3pt,
                minimum height=8mm, text width=2.15cm},
    enc/.style={box, fill=blue!8},
    head/.style={box, fill=orange!8},
    gtenc/.style={box, fill=red!8},
    dec/.style={box, fill=green!8},
    outbox/.style={box, fill=gray!10},
    loss/.style={box, fill=yellow!15, text width=6.6cm},
    arr/.style={thick, ->},
    das/.style={thick, dashed, ->},
]
% ---- training column x=0, vertical spacing 1.6cm --------------------
\node[enc]   (C)    at (0, 0)       {Sparse $C$\\\scriptsize (6 views)};
\node[enc]   (E)    at (0, -1.6)    {Encoder $E$\\\scriptsize UNet, 128 ch.};
\node[enc]   (TF)   at (0, -3.2)    {Geom. transformer\\\scriptsize anchor fusion};
\node[head]  (PSI)  at (0, -4.8)    {Pooled context\\$\psi(C)$};
\node[box]   (CONC) at (0, -6.4)    {$\psi(C)\oplus h_Y$};
\node[head]  (Q)    at (0, -8.0)    {Posterior head\\$q(z\mid C,Y)$};
\node[box]   (ZQ)   at (0, -9.6)    {$z\sim q(z\mid C,Y)$};
\node[dec]   (FILM) at (0, -11.2)   {FiLM modulation};
\node[dec]   (DEC)  at (0, -12.8)   {Decoder\\$g\circ f_{\mathrm{GS}}^{(z)}(C)$};
\node[outbox](XQ)   at (0, -14.4)   {$\hat{x}^{\,(q)}$};
\node[loss]  (LOSS) at (4.6, -15.3) {$\mathcal{L}=\mathcal{L}_{\mathrm{vol}}^{(q)}
   +\beta_t\,\mathrm{KL}[q\|p]$\\[1pt]
   $+\,w_{\mathrm{prior}}\,\mathcal{L}_{\mathrm{vol}}^{(p)}$};

\node[gtenc] (Y)    at (6.8, -0.35) {GT volume $Y$};
\node[gtenc] (GTE)  at (6.8, -1.95) {GTEncoder\\\scriptsize $1{\to}16{\to}16{\to}32$};
\node[gtenc] (hY)   at (6.8, -3.55) {$h_Y\in\mathbb{R}^{32}$};
\node[head]  (PriH) at (9.8, -4.8)  {Prior head\\$p(z\mid C)$};
\node[box]   (KL)   at (9.8, -9.6)  {$\mathrm{KL}[q\|p]$};

\draw[arr] (C) -- (E) -- (TF) -- (PSI);
\draw[arr] (Y) -- (GTE) -- (hY);
\draw[arr] (hY.south) |- (CONC.east);
\draw[arr] (PSI) -- (CONC) -- (Q) -- (ZQ) -- (FILM) -- (DEC) -- (XQ);
\draw[arr] (PSI.east) -| (PriH.west);
\draw[arr] (XQ.south) -| (LOSS.north west);
\draw[das] (Q.east) -| (KL.west);
\draw[das] (PriH.south) -| (KL.north);
\draw[arr] (KL.south) -| (LOSS.north east);

\node[draw=red!70!black, dashed, fit=(Y)(GTE)(hY), inner sep=4pt,
      label={[red!70!black, font=\scriptsize]above:training only (discarded at deployment)}] {};
\node[draw=black!60, dashed, fit=(C)(LOSS)(KL)(Y), inner sep=6pt,
      label={[font=\small\bfseries]above:Training}] {};

% ---- deployment row y=-17.6, spacing 2.1cm --------------------------
\node[enc, text width=1.4cm]   (DC)    at (0.0, -17.6)  {Sparse $C$\\\scriptsize (6 views)};
\node[enc, text width=1.4cm]   (DE)    at (2.1, -17.6)  {Enc. $E$\\\scriptsize + transformer};
\node[head, text width=1.4cm]  (DPSI)  at (4.2, -17.6)  {$\psi(C)$\\\scriptsize pooled\\context};
\node[head, text width=1.4cm]  (DPRI)  at (6.3, -17.6)  {Prior head\\$p(z\mid C)$};
\node[box, text width=1.4cm]   (DZ)    at (8.4, -17.6)  {$z_n\sim p(z\mid C)$\\\scriptsize $n{=}1,\dots,8$};
\node[dec, text width=1.4cm]   (DFILM) at (10.5, -17.6) {FiLM\\\scriptsize layers};
\node[dec, text width=1.4cm]   (DDEC)  at (12.6, -17.6) {Decoder\\\scriptsize context\\features};
\node[outbox, text width=1.4cm](DX)    at (14.7, -17.6) {$\hat{x}$\\[1pt]$\widehat{\mathrm{Var}}(x)$};

\draw[arr] (DC) -- (DE) -- (DPSI) -- (DPRI) -- (DZ) -- (DFILM) -- (DDEC) -- (DX);

\node[draw=teal!70!black, dashed, fit=(DC)(DX), inner sep=3pt,
      label={[font=\small\bfseries, teal!70!black]above:Deployment (no GT, no dense views)}] {};
\end{tikzpicture}}
\end{figure}